\section{Every odd exponent at least five}
\begin{theorem}[Source Theorem 1.4]\label{F2.Existence}
For every odd prime $p$ and every odd $m\geq5$, and for every even
$m\geq p+3$, there is a non-tight chiral polyhedron of order
$2p^m$. The following two constructions realize the exact source
presentations and have type $\{p^2,2p^e\}$, with $e$ specified below.
\end{theorem}
The proof is the complete construction in this section and the next.

First construct the underlying source $p$-group for every $n=m-1\geq2$,
without a parity restriction. Let $O=\mathbb Z[\zeta]$, $\pi=\zeta-1$,
where $\zeta$ is a primitive $p$th root, and write
$n=(e-1)(p-1)+r$, $1\leq r\leq p-1$. Set
\[
M=O/\pi^nO,\qquad s=1,\qquad Tv=\zeta v,
\qquad z=p^{e-1}\pi^{r-1}.
\]
The common lattice theorem gives $|M|=p^n$, $|s|=p^e$ and
$0\ne z\in\pi^{n-1}O/\pi^nO$, so $|z|=p$ and $Tz=z$.
The cyclic-extension construction in the proof of the carry theorem,
which needs only $T^p=1$ and $Tz=z$ to construct $P$, gives
\[
P=(M\rtimes\mathbb Z)/\gen{(-z,p)},\qquad
|P|=p^{n+1},\qquad v^\beta=Tv,\quad\beta^p=z.
\]
The powers of $\zeta=1+\pi$ span $M$, so $P=\gen{\beta,s}$.
As $M$ is abelian, $[P,P]=\pi M$ and inductively
$\gamma_j(P)=\pi^{j-1}M$ for every $j\geq2$.
Thus $\gamma_n(P)\ne1$ and $\gamma_{n+1}(P)=1$: its class is $n$.
The norm $\sum_{i=0}^{p-1}T^i$ is zero, so
$(\beta s)^p=\beta^p=z$ and $|\beta s|=p^2$.
More generally $(\beta^kv)^p=kz$ for $1\leq k<p$ and $v\in M$,
since multiplication by $k$ permutes the powers of $T$.
Every element outside $M$ therefore has order $p^2$.
An automorphism fixing $\beta$ and inverting $s$ would invert all its
$\beta$-conjugates, hence all of $M$, while fixing $z=\beta^p$.
This contradicts $|z|=p>2$. These conclusions require no reflection $S$.

For the literal source presentation, put $s_i=\pi^{i-1}$.
The mixed coordinate orders are precisely those printed in Section 3;
$[s_i,\beta]=s_{i+1}$ and
\[
[s_{p-1},\beta]=\pi^{p-1}
=-\sum_{i=1}^{p-1}\binom pi\pi^{i-1}.
\]
The carry is $\beta^p=s_r^{p^{e-1}}$.
Conversely those relations collect every source word into
$\beta^j\prod s_i^{a_i}$ with $0\leq j<p$ and the indicated coordinate
ranges, including coordinates of order one. They bound the presented
order by $p^{n+1}$, and the constructed model realizes equality.
This proves the exact presentation for every $n\geq2$; it does not
assert the classification of all maximal-class groups.

For the chiral construction now assume that $n\geq4$ is even.
Then $r$ is even and $r\geq2$. Put $Sv=\bar v$.
Conjugation satisfies $\bar\pi=-\zeta^{-1}\pi$, hence
$Sz=(-1)^{r-1}z=-z$, and $STS=T^{-1}$ and $Ss=s$.
In the direct coordinates $1,\pi,\ldots,\pi^{p-2}$,
$z$ is a nonzero element of the $r$th coordinate whereas $s$ is
in the first; thus $\gen z\cap\gen s=0$.
The carry theorem of~\cite{UF} gives order $2p^{n+1}$ and type
$\{p^2,2p^e\}$. This is non-tight since $n+1>e+2$ for even
$n\geq4$ and $p\geq3$.

The two source formulas for
conjugation are UF.Conjugation, which also explains the long
binomial verification without retaining it as an external proof premise.

